%%%  Ukázkový text a dokumentace stylu pro text závěrečné (bakalářské a diplomové) práce na KI PřF UP v Olomouci
%%%  Copyright (C) 2026 Jesse Sadowý, <jesse.sadowy23@gmail.com>
%%%  Copyright (C) 2026 Jan Konečný, <jan.konecny@upol.cz>


%%  Pro získání PDF souboru dokumentu je třeba tento zdrojový text v
%%  LaTeXu přeložit (dvakrát) programem pdfLaTeX.

%%  V případě použití programu BibLaTeX pro tvorbu seznamu literatury
%%  je poté ještě třeba spustit program Biber s parametrem jméno
%%  souboru zdrojového textu bez přípony a následně opět (dvakrát)
%%  přeložit zdrojový text programem pdfLaTeX.

%%  Postup získání Postscriptového souboru je popsán v dokumentaci.


%%  Třída dokumentu implementující styl pro závěrečnou práci. Vybrané
%%  nepovinné parametry (ostatní v dokumentaci):

%%  'master' pro sazbu diplomové práce, jinak se sází bakalářská práce

%%  'program=kód' pro Váš studijní program/obor (specializaci), kódy
%%  pro diplomovou práci 'infoi' pro Informatiku (Obecná informatika),
%%  'infui' pro Informatiku (Umělá inteligence), 'ainfpst' pro
%%  Aplikovanou informatiku (Počítačové systémy a technologie), 'uinf'
%%  pro Učitelství informatiky pro střední školy, 'binf' pro
%%  Bioinformatiku, 'inf' pro Informatiku (bez specializací) a 'ainf'
%%  pro Aplikovanou informatiku (bez specializací), jinak je výchozí
%%  ainfvs pro Aplikovanou informatiku (Vývoj software), a pro
%%  bakalářskou práci 'infoi' pro Informatiku (Obecná informatika),
%%  'itp' pro Informační technologie v prezenční formě, 'itk' pro
%%  Informační technologie v kombinované formě, 'infv' pro Informatiku
%%  pro vzdělávání, 'binf' pro Bioinfomatiku, 'inf' pro Informatiku
%%  (bez specializací), 'ainfp' pro Aplikovanou informatiku (bez
%%  specializací) v prezenční formě, 'ainfk' pro Aplikovanou
%%  informatiku (bez specializací) v kombinované formě, jinak je
%%  výchozí infpvs pro Informatiku (Programování a vývoj software)

%%  'printversion' pro sazbu verze pro tisk (nebarevné logo a odkazy,
%%  odkazy s uvedením adresy za odkazem, ne odkazy do rejstříku),
%%  jinak verze pro prohlížeč

%%  'biblatex' pro zapnutí podpory pro sazbu bibliografie pomocí
%%  BibLaTeXu, jinak je výchozí sazba v prostředí thebibliography

%%  'language=jazyk' pro jazyk práce, jazyky english pro anglický,
%%  slovak pro slovenský, jinak je výchozí czech pro český

%%  'font=sans' pro bezpatkový font (Iwona Light), jinak je výchozí
%%  serif pro patkový (Latin Modern)

%%  'figures, tables, theorems a sourcecodes' pro sazbu seznamu
%%  obrázků, tabulek, vět a zdrojových kódů, jinak při =false se
%%  nesází (u theorems a sourcecodes výchozí)

\documentclass[
%  master,
  program=itk,
%  printversion,
  biblatex,
%  language=english,
%  font=sans,
  figures=false,
%  tables=false,
%  theorems,
%  sourcecodes,
  glossaries,
  index
]{kidiplom}

%% Informace pro úvodní strany. V jazyku práce (pokud není v komentáři
%% uvedeno česky) a anglicky. Uveďte všechny, u kterých není v
%% komentáři uvedeno, že jsou volitelné. Při neuvedení se použijí
%% výchozí texty. Text pro jiný než nastavený jazyk práce (nepovinným
%% parametrem language makra \documentclass, výchozí český) se zadává
%% použitím makra s uvedením jazyka jako nepovinného parametru.

%% Název práce, česky a anglicky. Měl by se vysázet na jeden řádek.
\title{Procedurální generování dungeonů ve 2D hrách}
\title[english]{Procedural Generation of Dungeons in 2D Games}

%% Volitelný podnázev práce, česky a anglicky. Měl by se vysázet na
%% jeden řádek. Výchozí je prázdný.
\subtitle{Implementace a srovnání algoritmů v herním enginu Godot}
\subtitle[english]{Implementation and Comparison of Algorithms in Godot Game Engine}

%% Jméno autora práce. Makro nemá nepovinný parametr pro uvedení
%% jazyka.
\author{Jesse Sadowý}

%% Jméno vedoucího práce (včetně titulů). Makro nemá nepovinný
%% parametr pro uvedení jazyka.
\supervisor{doc. RNDr. Jan Konečný, Ph.D.}

%% Volitelný rok odevzdání práce. Výchozí je aktuální (kalendářní)
%% rok. Makro nemá nepovinný parametr pro uvedení jazyka.
\yearofsubmit{2026}

%% Anotace práce, včetně anglické (obvykle překlad z jazyka
%% práce). Jeden odstavec!
\annotation{Bakalářská práce se zabývá procedurálním generováním dungeonů ve 2D hrách. Zkoumá různé metody a techniky používané pro generování a následně porovnává jejich efektivitu a výkonnost v prostředí herního enginu Godot.}

\annotation[english]{Bachelor's thesis focuses on procedural generation of dungeons in 2D games. It explores various methods and techniques used for generation and subsequently compares their efficiency and performance within the Godot game engine environment.}

%% Klíčová slova práce, včetně anglických. Oddělená (obvykle) středníkem.
\keywords{procedurální generování; dungeon; 2D hry; Godot; herní engine, algoritmy}
\keywords[english]{procedural generation; dungeon; 2D games; Godot; game engine, algorithms}
%% Volitelná specifikace příloh textu práce, i anglicky. Výchozí je
%% 'elektronická data v systému katedry informatiky / electronic data
%% in system of department of computer science'.
%\supplements{nejlepší software všech dob}
\supplements{test}
%\supplements[english]{the best software of all times}

%% Volitelné poděkování. Stručné! Výchozí je prázdné. Makro nemá
%% nepovinný parametr pro uvedení jazyka.
\thanks{Děkuji, děkuji, děkuji.}

%% Nastavení biblatex pro odstranění varování
\ExecuteBibliographyOptions{seconds=true}

%% Cesta k souboru s bibliografií pro její sazbu pomocí BibLaTeXu
%% (zvolenou nepovinným parametrem biblatex makra
%% \documentclass). Použijte pouze při této sazbě, ne při (výchozí)
%% sazbě v prostředí thebibliography.
\bibliography{bibliografie.bib}

%% Další dodatečné styly (balíky) potřebné pro sazbu vlastního textu
%% práce.
\usepackage{lipsum}
\usepackage{longtable}

\begin{document}
%% Sazba úvodních stran -- titulní, s bibliografickými údaji, s
%% anotací a klíčovými slovy, s poděkováním a prohlášením, s obsahem a
%% se seznamy obrázků, tabulek, vět a zdrojových kódů (pokud jejich
%% sazba není vypnutá).
\maketitle

%% Vlastní text závěrečné práce. Pro povinné závěry, před přílohami,
%% použijte prostředí kiconclusions. Povinná je i příloha s obsahem
%% elektronických dat.

%% -------------------------------------------------------------------

\newcommand{\BibLaTeX}{\textsc{Bib}\LaTeX}

% \noindent\textcolor{red}{\LARGE Upozornění: Následující text
%   dokumentace stylu, vyjma přílohy~\ref{sec:ObsahData}, je rozpracovaná
%   a (značně) neúplná verze!!!}

\section{Úvod}

\section{Procedurální generování obsahu ve hrách}

\subsection{Vymezení a definice pojmu}

Procedurální generování obsahu ve hrách lze vymezit jakožto algoritmickou tvorbu herního obsahu s limitovaným nebo nepřímým zásahem hráče. Tento přístup nám umožňuje automatizovanou tvorbu různých prvků hry bez nutnosti jejich úpravy. Herní obsah zahrnuje široké spektrum prvků jako jsou mapy, úrovně, pravidla, úkoly, zbraně či předměty.

PCG se tedy neomezuje pouze na herní prostřední, ale také i na výše zmíněné prvky.

V současné době je pojem generování obsahu často spojován s využitím metod AI, avšak zde se pojednává výhradně o algoritmické metody, založené na předem definovaných pravidlech a postupech.

Nedílným klíčovým pojmem jsou hry, který je však těžké definovat a je nutno zmínit, že v tomto případě se jedná o video hry především pro dnešní počítače.

Termíny "procedurální" a "generování" zdůrazňují, že se zabýváme počítačovými procedurami a algoritmy, které nám generují obsah. PCG spustí člověk na počítači a "něco" nám vygeneruje.

\subsection{Motivace a důvody využití procedurálního generování}

Hlavním důvodem pro použití PCG je, že obsah je v podstatě generován sám, takže odstraňuje nutnost kreslit každou novou mapu či položku zvlášť, vytvrářet logiku pro další úrovně, případně v každé truhle, kterou ve hře najdeme, specifikovat, co tam bude. Nejenže tím zjedodušíme tvorbu, ale výrazně ušetříme čas i náklady na vývoj dané hry.

Dalším výrazným přínosem je to, že PCG přidává momenty překvapení a znovuhratelonsti, jelikož obsah se bude při každém průchodu lišit a hráč tím pádem nikdy neví, co se stane a jaká výzva ho čeká. Dále nám to umožňuje hru přizpůsobit samotnému hráči podle toho, jak hraje. Příkladem může být proměnlivá obtížnost, která zůstane vyvážená a hra tím nebude příliš jednoduchá či složitá.
Vhodné je použít PCG i při umisťování herních předmětů či nepřátel, tím pádem bude každý průchod jedinečný a opět to dodá hře unikátnost a znovuhratelnost.

Na první pohled to tak nemusí vypadat, ale tento způsob dodává hře i více kreativních prvků, jelikož daný algoritmus může přinést výsledky, na které by vývojář nemusel nikdy přijít. Člověk má tendenci používat to, co už někdy viděl a tímto bude obsah tvořen zcela novým a nečekaným způsobem.

Nakonec je vhodné podotknout, že díky PCG se naučíme porozumět i designu her samotných. Jak uvádí Shaker et al. \cite{pcgbook}, proces tvorby generátoru vyžaduje hluboké porozumění herním mechanikám a principům dobrého designu, které musíme převést do algoritmické podoby.

\subsection{Nevýhody procedurálního přístupu}
\subsection{Využití v herním průmyslu}

\section{Dungeony ve 2D hrách}

\subsection{Charakteristické vlastnosti dungeonových map}
\subsection{Typické herní prvky dungeonů}
\subsection{Požadavky na generování dungeonových map}

\subsection{Přehled vybraných metod generování dungeonů}

\subsubsection{Generování pomocí náhodného umístění místností}
\subsubsection{Rozdělování prostoru metodou Binary Space Partitioning}
\subsubsection{Generování pomocí buněčných automatů}
\subsubsection{Metoda náhodné procházky}

\section{Použité technologie}
\subsection{Herní engine Godot}
\subsection{Skriptovací jazyk GDScript}
\subsection{Reprezentace dlaždicové mapy}
\subsection{Reprezentace hráče pro vizualizaci}

\section{Návrh systému generování}

\subsection{Datová reprezentace dungeonové mapy}
\subsection{Architektura generátoru}
\subsection{Společné kroky generování}
\subsection{Možnosti rozšíření systému}

\section{Implementace generovacích metod}

\subsection{Generování pomocí místností a chodeb}
\subsubsection{Princip metody}
\subsubsection{Generování místností}
\subsubsection{Propojování chodeb}
\subsubsection{Ukázky výsledných map}

\subsection{Generování pomocí bludiště}
\subsubsection{Princip metody}
\subsubsection{Postup generování}
\subsubsection{Ukázky výsledných map}

\subsection{Binary Space Partitioning}
\subsubsection{Princip rozdělování prostoru}
\subsubsection{Postup generování}
\subsubsection{Ukázky výsledných map}

\subsection{Metoda náhodné procházky (Drunkard's Walk)}
\subsubsection{Princip metody}
\subsubsection{Postup generování}
\subsubsection{Ukázky výsledných map}

\subsection{Další možné metody generování}
\subsubsection{Stručný přehled}
\subsubsection{Možnosti budoucí implementace}

\section{Vyhodnocení a porovnání metod}

\subsection{Zvolená hodnoticí kritéria}
\subsection{Porovnání výsledných map}
\subsection{Výpočetní náročnost metod}
\subsection{Shrnutí výsledků}

\section{Diskuze}

\subsection{Zkušenosti z implementace}
\subsection{Silné a slabé stránky jednotlivých metod}
\subsection{Možnosti dalšího rozšíření práce}

% TODO doplnit zkratky
%% Závěry práce
\begin{kiconclusions}
% TODO: Doplnit český závěr
\end{kiconclusions}

\begin{kiconclusions}[english]
% TODO: Doplnit anglický závěr
\end{kiconclusions}

%% Přílohy
\appendix

\section{Ukázky zdrojového kódu}
% TODO: Doplnit ukázky kódu

\section{Odkaz na repozitář projektu}
% TODO: Doplnit odkaz na repozitář

\section{Zkratky}
PGC - Procedural Content Generation
% TODO doplnit zkratky

%% Obsah elektronických dat.
\section{Obsah elektronických dat} \label{sec:ObsahData}

Na samotném konci textu práce je uveden stručný popis obsahu
elektronických dat odevzdaných v systému katedry informatiky spolu s
textem. Tato data jsou nedílnou součástí práce a tvoří (datovou)
přílohu textu práce. Povinné položky struktury dat jsou:

\begin{description}

\item[\texttt{text/}] \hfill \\
  Adresář s textem práce ve formátu PDF, vytvořený s~použitím
  závazného stylu KI PřF UP v~Olomouci pro závěrečné práce, včetně
  všech (textových) příloh, a~všechny soubory potřebné pro
  bezproblémové vytvoření PDF dokumentu textu (případně v~ZIP
  archivu), tj.~zdrojový text textu a příloh, vložené obrázky, apod.

\item[\texttt{README.*}] \hfill \\
  Textový soubor (s příponou např. \texttt{.txt}) s informacemi o
  opakovatelném způsobu použití ostatních dat práce -- typicky plně
  reprodukovatelný co nejúplnější funkční postup zprovoznění software
  vytvořeného v~rámci práce, tzn. jeho případné instalace/nasazení a
  spuštění, včetně uvedení všech požadavků pro bezproblémový provoz;
  za zprovoznění software se nepovažuje zpřístupnění (např. po
  Internetu) již někde zprovozněného software.

\item[\texttt{*}] \hfill \\
  Adresáře a soubory s veškerými ostatními autorskými daty práce
  (případně v~ZIP archivu) -- typicky spustitelné a další soubory
  software vytvořeného v rámci práce potřebné pro bezproblémový provoz
  software, případně jeho instalační program, a kompletní zdrojové
  texty software a další data nutná pro plně reprodukovatelné korektní
  vytvoření spustitelných souborů.

\end{description}

Dále mohou data obsahovat například:

\begin{itemize}

%\item[\texttt{data/}] \hfill \\
\item
  ukázková a~testovací data použitá v~práci nebo pro potřeby posouzení
  práce v rámci její obhajoby,

%\item[\texttt{literature/}] \hfill \\
\item
  položky bibliografie v elektronické podobě, příp.~jiná relevantní literatura
  a dokumentace vztahující se k~práci,

%\item[\texttt{install/}] \hfill \\
\item
  cizí data (software) potřebná pro bezproblémové použití autorských
  dat práce (software), která nejsou standardní součástí
  předpokládaného (softwarového) vybavení uživatele.

\end{itemize}

U~veškerých cizích obsažených materiálů jejich
zahrnutí dovolují podmínky pro jejich veřejné šíření nebo přiložený souhlas
držitele práv k užití. Pro všechny použité (a~citované) materiály,
u~kterých toto není splněno a~nejsou tak obsaženy, je uveden
jejich zdroj, např.~webová adresa, v~bibliografii nebo textu práce
nebo souboru \texttt{README.*}.

%% -------------------------------------------------------------------

%% Sazba volitelného seznamu zkratek, za přílohami.
\printglossary

%% Sazba povinné bibliografie, za přílohami (případně i za seznamem
%% zkratek). Při použití BibLaTeXu použijte makro
%% \printbibliography. jinak prostředí thebibliography. Ne obojí!

%% Sazba i v textu necitovaných zdrojů, při použití
%% BibLaTeXu. Volitelné.
\nocite{*}
%% Vlastní sazba bibliografie při použití BibLaTeXu.
\printbibliography

%% Bibliografie, včetně sazby, při NEpoužití BibLaTeXu.
% \begin{thebibliography}{9}
%\bibitem{kniha2} \uppercase{Hawke}, Paul. NanoHttpd: Light-weight HTTP server designed for embedding in other applications. GitHub [online]. 2014-05-12. [cit. 2014-12-06]. Dostupné z: \url{https://github.com/NanoHttpd/nanohttpd}
%
%\bibitem{jeske13} \uppercase{Jeske}, David; \uppercase{Novák}, Josef. Simple HTTP Server in \csharp: Threaded synchronous HTTP Server abstract class, to respond to HTTP requests. CodeProject: For those who code [online]. 2014-05-24. [cit. 2014-12-06]. Dostupné z: \url{http://www.codeproject.com/Articles/137979/Simple-HTTP-Server-in-C}
%
%\bibitem{uzis2012} \uppercase{ÚSTAV ZDRAVOTNICKÝCH INFORMACÍ A STATISTIKY ČR}. Lékaři, zubní lékaři a farmaceuti 2012 [online]. Praha 2, Palackého náměstí 4: Ústav zdravotnických informací a statistiky ČR, 2012 [cit. 2014-12-06]. ISBN 978-80-7472-089-5. Dostupné z: \url{http://www.uzis.cz/publikace/lekari-zubni-lekari-farmaceuti-2012}
% \end{thebibliography}

%% Sazba volitelného rejstříku, za bibliografií.
\printindex

\end{document}

%%% Local Variables:
%%% mode: latex
%%% TeX-master: t
%%% End:
